% TOP_PROBLEM: {"id": 34, "title": "The Volume Conjecture", "category_id": 7, "category": {"id": 7, "name": "topology", "display_name": "Topology", "description": "Properties preserved under continuous deformations.", "slug": "topology", "order_index": 7, "created_at": "2026-07-31T15:26:25.671Z"}, "status": "open"}
% FIRST_POSTED: 2026-09-27
% !TeX program = lualatex
% Compile twice: lualatex 34.tex
% All references and prose are embedded; no BibTeX database or external inputs.
\documentclass[11pt,a4paper]{article}
\usepackage[margin=24mm,headheight=14pt]{geometry}
\usepackage{fontspec}
\setmainfont{Latin Modern Roman}
\setsansfont{Latin Modern Sans}
\setmonofont{Latin Modern Mono}
\usepackage{amsmath,amssymb,mathtools}
\usepackage{unicode-math}
\setmathfont{Latin Modern Math}
\usepackage{microtype}
\usepackage{enumitem,xurl,needspace}
\usepackage[dvipsnames]{xcolor}
\usepackage{hyperref}
\hypersetup{unicode=true,colorlinks=true,linkcolor=MidnightBlue,urlcolor=MidnightBlue,
 pdftitle={500 Mathematical Problem Statements},pdfauthor={Source-annotated compilation},bookmarksnumbered=true}
\usepackage{bookmark}
\usepackage{tocloft}
\setlength{\cftsecnumwidth}{3em}
\cftsetpnumwidth{2.5em}
\cftsetrmarg{4em}
\usepackage{titlesec}
\titleformat{\section}{\Large\bfseries\raggedright}{\thesection}{0.7em}{}
\usepackage{fancyhdr}
\pagestyle{fancy}\fancyhf{}
\fancyhead[L]{\small\sffamily Mathematical problem statements}
\fancyhead[R]{\small Source edition 22}
\fancyfoot[C]{\thepage}
\setlength{\parindent}{0pt}
\setlength{\parskip}{5pt plus 1pt}
\setlength{\emergencystretch}{3em}
\setcounter{tocdepth}{1}
\setcounter{secnumdepth}{1}
\setlist[itemize]{leftmargin=3em,itemsep=5pt,topsep=4pt,parsep=0pt}
\allowdisplaybreaks
\newcommand{\N}{\mathbb{N}}
\newcommand{\Z}{\mathbb{Z}}
\newcommand{\Q}{\mathbb{Q}}
\newcommand{\R}{\mathbb{R}}
\newcommand{\C}{\mathbb{C}}
\newcommand{\F}{\mathbb{F}}
\newcommand{\A}{\mathbb{A}}
\newcommand{\PP}{\mathbb P}
\newcommand{\E}{\mathbb{E}}
\newcommand{\eps}{\varepsilon}
\newcommand{\dd}{\,\mathrm d}
\newcommand{\sref}[2]{\hyperlink{source:#1:#2}{[S#2]}}
\newcommand{\eref}[2]{\hyperlink{extra:#1:#2}{[E#2]}}
% Turn the next switch off for a shorter reading copy. The quotations preserve
% every exactTarget field, including qualifications omitted from a short summary.
\newif\ifshowcatalogue
\showcataloguetrue
\newcommand{\cataloguescope}[1]{\ifshowcatalogue
 \paragraph{Exact scope in the supplied catalogue.}
 \begin{quote}\small #1\end{quote}\fi}
\hypersetup{pdftitle={Problem 34 - Knot volume conjecture},pdfauthor={Open Problems Math research notebook}}
\fancyhead[L]{\small\sffamily Open Problems Math \textbar\ Research notebook}
\fancyhead[R]{\small\sffamily 121 \textbar\ 27 September 2026}
\numberwithin{equation}{section}
\begin{document}
\setcounter{section}{0}
% ================================================================
% problemId: 34
% Canonical problem ID: 34
% exactTarget SHA-256: c8c5773f55ad967e9e98f3555dbc6cb836b858c1b085ddc39f12c6e621048aa7
\clearpage
\section*{\texorpdfstring{The Volume Conjecture}{The Volume Conjecture}}\label{34}
\subsection{Definitions and mathematical statement}
For a knot $K\subset S^3$, let $J_{K,N}(q)$ be the $N$-colored Jones polynomial normalized so that the unknot has value one. A hyperbolic knot has a complete finite-volume hyperbolic metric on its complement. The conjecture is
\[
 \lim_{N\to\infty}\frac{2\pi}{N}\log\left|J_{K,N}(e^{2\pi i/N})\right|
 =\operatorname{Vol}(S^3\setminus K).
\]
Both the normalization and the root-of-unity convention matter. The statement includes existence of the limit for every hyperbolic knot; a numerical fit or a result for one family is not the universal conclusion.
\par\smallskip\textit{Formulation and concept references:} \sref{34}{1}.
\cataloguescope{Prove or disprove that for every hyperbolic knot K in S\textasciicircum{}3, the limit as N tends to infinity of (2*pi/N) log|J\_\{K,N\}(exp(2*pi*i/N))| equals the hyperbolic volume of S\textasciicircum{}3 minus K, using the standard normalized N-colored Jones polynomial.}
\subsection{Short English statement}
Does the exponential growth of colored Jones values at roots of unity recover the hyperbolic volume of every hyperbolic knot complement?
% BEGIN RESEARCH ADDITION 121
\subsection{Research attempt}

\paragraph{Current frontier.}
The formulation source studies quantum modular forms and, in particular,
the figure-eight knot.  Its introduction gives a positive product-sum
formula for the figure-eight Kashaev invariant and the matching volume
normalization.  It does not prove the stated limit for every hyperbolic
knot. \sref{34}{1}
The distinction is analytically important: positivity makes the leading
exponential rate accessible without proving cancellation estimates.
The attempt below rederives that rate, including endpoint control, then
isolates what a general complex state sum would additionally require.

\paragraph{Attack routes.}
One can try a saddle-point expansion of a knot state sum, deformation
of integration contours toward a geometric critical point, or a
positive-sum comparison.  The first needs a uniform remainder and control
of all other saddles; the second needs a justified contour deformation
with no omitted residues; the third needs a lower bound preventing
exponential phase cancellation.  We first make the positive case fully
rigorous, then formulate the missing phase condition without treating it
as an automatic consequence of the existence of a geometric saddle.

\paragraph{Attempt I: a uniform Laplace argument for the figure-eight knot.}
For the normalization in the statement, the figure-eight formula is
\begin{equation}\label{34:sum}
 J_{4_1,N}(e^{2\pi i/N})=
 \sum_{k=0}^{N-1}\prod_{j=1}^{k}4\sin^2(\pi j/N).
\end{equation}
The empty product is one.  The source writes the equivalent invariant
using its root-of-unity convention and explicitly gives the positive
formula. \sref{34}{1}, introduction, equations (1)--(3).
Put $f(x)=\log(2\sin\pi x)$ on $(0,1)$ and
$F(x)=2\int_0^x f(t)\,dt$.  The endpoint singularities are integrable.
We claim the uniform estimate
\begin{equation}\label{34:riemann}
 \max_{0\leq k<N}\left|
 2\sum_{j=1}^k f(j/N)-NF(k/N)\right|=O(\log N).
\end{equation}
To prove it, write
\[
 f(x)=\log x+\log(1-x)+g(x),\qquad
 g(x)=\log\frac{2\sin\pi x}{x(1-x)}.
\]
The function $g$ extends smoothly to $[0,1]$, so its partial Riemann-sum
error is uniformly $O(1)$.  For $\log x$, monotone integral comparison,
or $\sum_{j\leq k}\log j=\log(k!)$ and Stirling's inequalities, bounds
the error by $O(\log N)$ uniformly in $1\leq k<N$.  Applying the same
comparison to the reversed list $N-j$ handles $\log(1-x)$, including
$k=N-1$.  At $k=0$ both expressions vanish.  Adding proves
\eqref{34:riemann}; no unjustified uniform convergence near $0$ or $1$
is being used.

Since $F'(x)=2\log(2\sin\pi x)$, it decreases on $(0,1/6)$,
increases on $(1/6,5/6)$, and decreases on $(5/6,1)$.
The identity $\int_0^1\log(2\sin\pi x)\,dx=0$ gives
$F(0)=F(1)=0$.  Moreover $F(5/6)>0$, by reflection of the negative
integral on $(0,1/6)$.  Hence its global maximum is $F(5/6)$.
This maximum is attained in the interior, where $F$ is smooth; choosing
a nearest grid point changes $NF$ by $O(1)$ (in fact by $O(N^{-1})$).
If $a_{N,k}$ is the $k$th positive summand in \eqref{34:sum}, then
\[
 \max_k a_{N,k}\leq\sum_k a_{N,k}\leq N\max_k a_{N,k}.
\]
Together with \eqref{34:riemann}, this proves
\begin{equation}\label{34:rate}
 \log J_{4_1,N}(e^{2\pi i/N})
 =2N\int_0^{5/6}\log(2\sin\pi x)\,dx+O(\log N).
\end{equation}
The standard ideal-tetrahedron volume expression, also specified in the
source, is
\[
 \operatorname{Vol}(S^3\setminus4_1)
 =4\pi\int_0^{5/6}\log(2\sin\pi x)\,dx.
\]
Thus \eqref{34:rate} proves the requested normalization for this knot,
with error $O(\log N/N)$ after scaling.  This is a reconstruction of a
known solved case, not a new universal result. \sref{34}{1}

\paragraph{Attempt II: formulate the precise cancellation gap.}
Consider a general finite state sum $Z_N=\sum_{s\in\mathcal S_N}z_{N,s}$.
Suppose one has proved
\[
 \sum_s|z_{N,s}|\leq\exp(N(v+o(1))).
\]
This gives an upper bound on $N^{-1}\log|Z_N|$, but not a lower bound.
Here is one sufficient condition that would repair the argument.  Find
subsets $\mathcal A_N$, phases $\theta_N$, and numbers $c_N>0$ with
$\log c_N=o(N)$ such that
\begin{align}\label{34:phase}
 \operatorname{Re}(e^{-i\theta_N}z_{N,s})&\geq c_N|z_{N,s}|
                 &&(s\in\mathcal A_N),\\
 W_N:=\sum_{s\in\mathcal A_N}|z_{N,s}|&\geq\exp(N(v-o(1))),\notag\\
 \sum_{s\notin\mathcal A_N}|z_{N,s}|&\leq\tfrac12c_NW_N.\notag
\end{align}
Then
\[
 |Z_N|\geq\operatorname{Re}(e^{-i\theta_N}Z_N)
       \geq\tfrac12c_NW_N,
\]
and the lower and upper exponential rates agree.  This elementary
phase-sector lemma gives an exact target for a contour or state-sum
construction: not just a large critical value, but a phase-coherent
contribution dominating the rest at the exponential scale.

For a knot $K$, this route would require $v=\operatorname{Vol}(S^3
\setminus K)/(2\pi)$ and a state sum matching the \emph{normalized}
colored Jones value at every sufficiently large integer $N$.  Neither
that phase package nor an equivalent noncancellation theorem is proved
here for arbitrary $K$.  Polynomially many terms alone does not help.
For example, the two terms
\[
 z_{N,1}=e^{Nv},\qquad z_{N,2}=-e^{Nv}+e^{N(v-c)},\quad c>0,
\]
have absolute exponential rate $v$ and sum of rate $v-c$.
Replacing the last term by zero makes the sum vanish.  Therefore even
exact identification of the maximal real part of an action cannot
justify the lower bound or the existence of the logarithmic limit.

\paragraph{Stress test.}
Equation \eqref{34:sum} is special: all its terms are positive, and none
uses the singular factor at $j=N$.  The Riemann estimate explicitly
handles terms adjacent to both singular endpoints.  Taking an absolute
value term by term is harmless for its upper bound but is not promoted
to a lower bound for other knots.  The phase-sector condition includes
a bound on the omitted states, so it cannot hide a cancelling saddle
outside the chosen neighborhood.

The companion computation evaluates \eqref{34:sum} in logarithmic
form for a finite sequence of $N$, using high-precision arithmetic and
log-sum-exp to avoid overflow.  It also checks the cancellation example.
The displayed error estimate is proved analytically; the decimal
values are diagnostics, not rigorous interval enclosures or a proof
for an additional knot.  No claim about all hyperbolic knots follows
from these finite computations.

\paragraph{Outcome.}
\textbf{Outcome: Exploratory approach with unresolved gap.}

The positive-sum route has been carried through with a uniform endpoint
estimate for the figure-eight knot.  A phase-sector lemma identifies a
sufficient lower-bound mechanism for more general state sums, and an
explicit example disproves a shortcut based only on their largest
absolute terms.  The universal noncancellation and geometric matching
steps remain unproved.  No new knot case, counterexample, or priority
claim is asserted.
% END RESEARCH ADDITION 121
\subsection{Sources}
\begin{itemize}\raggedright
\item[\textnormal{[S1]}]\hypertarget{source:34:1}{} C. Aistleitner and B. Borda, Journal of the European Mathematical Society 27 (2025), 5043-5092.
\url{https://ems.press/content/serial-article-files/51448?nt=1}.
{\small\textit{Catalogue formulation source.}}
\end{itemize}


\end{document}
